\section{Weather, caution, input design and robustness to pose error}\label{sec:weather}

\begin{plainwords}
This section answers the original question: does the driving get worse in fog and rain? The answer is no, in sparse traffic, for the perception damage we measured. It also reports three side studies: a driver that becomes cautious when its perception is unreliable, which input design works best, and what happens when the car's position estimate on the map is wrong.
\end{plainwords}
\begin{cautionbox}[Compared with GitHub]
The GitHub README lists the weather evaluation as ``preliminary tooling'' and states that it ``does not change vehicle dynamics or road friction'' and is ``not evidence of physical fog/rain robustness or adaptive driving''. The results below complete that picture within the same limits: weather acts only through the camera pipeline and, in this report, only through a measured degradation of what the model sees.
\end{cautionbox}


This section contains the results that bear on the original question (``can the policy keep driving, and drive more economically, when its perception degrades?''), in the order in which the evidence was produced.

\subsection{Detector-dropout proxy: no effect (pilot, 40 episodes)}
With the dropout proxy (each vehicle removed from the grid with the measured retention probability of the weather), three models were evaluated on seeds 1010--1049 (Table~\ref{tab:dropout}). The success rates are \emph{identical} across weather conditions for the unmodified-grid student and differ by at most one episode (0.025) for the other two models. Even at rain 1.0, where only 4\% of the vehicles survive, nothing changes.
\begin{table}[H]\centering\small
\caption{Success under the detector-dropout proxy, sparse traffic, seeds 1010--1049 (pilot; 40 episodes per cell). ``Share of vehicles kept'' is the measured detector retention used by the proxy.}\label{tab:dropout}
\resizebox{\textwidth}{!}{\begin{tabular}{@{}llrrr@{}}\toprule
Model & Weather & Share of vehicles kept & Success (40 episodes) & 95\% interval\\\midrule
Baseline behaviour cloning, sparse-traffic demonstrations & clear & 1.00 & 0.775 & [0.62, 0.88]\\
 & fog 1.0 & 0.27 & 0.800 & [0.65, 0.90]\\
 & rain 0.6 & 0.16 & 0.800 & [0.65, 0.90]\\
 & rain 1.0 & 0.04 & 0.800 & [0.65, 0.90]\\
Baseline behaviour cloning, no-traffic demonstrations & clear & 1.00 & 0.875 & [0.74, 0.95]\\
 & fog 1.0 & 0.27 & 0.875 & [0.74, 0.95]\\
 & rain 0.6 & 0.16 & 0.875 & [0.74, 0.95]\\
 & rain 1.0 & 0.04 & 0.875 & [0.74, 0.95]\\
Dataset-aggregation student, unmodified view & clear & 1.00 & 0.875 & [0.74, 0.95]\\
 & fog 1.0 & 0.27 & 0.875 & [0.74, 0.95]\\
 & rain 0.6 & 0.16 & 0.875 & [0.74, 0.95]\\
 & rain 1.0 & 0.04 & 0.875 & [0.74, 0.95]\\
Headline model (map-overlay view) & clear & 1.00 & 0.925 & [0.80, 0.97]\\
 & fog 1.0 & 0.27 & 0.925 & [0.80, 0.97]\\
 & rain 0.6 & 0.16 & 0.925 & [0.80, 0.97]\\
 & rain 1.0 & 0.04 & 0.950 & [0.83, 0.99]\\
\bottomrule\end{tabular}
}
\end{table}

\subsection{Calibrated degradation: no effect (reported, 200 episodes)}
After the measurement of Section~\ref{sec:perception} the environment was rebuilt so that vehicles can both disappear and be displaced by a wrong stereo range (\pth{bev_rl/adaptive/degraded_env.py}); weather conditions: clear, fog 1.0, rain 0.6, rain 1.0.
Following the pre-registration, the first test (H1) asked whether a clean-trained model gets \emph{worse} under the calibrated degradation: proposed model, 100 paired episodes, seeds 5000--5099: clear 0.89, fog 1.0: 0.90, rain 0.6: 0.89, rain 1.0: 0.89 (exact McNemar $p=1.0$ for every degraded cell against clear). \textbf{H1 failed}: the pre-registered consequence is that adaptation cannot be demonstrated on success in sparse traffic, and everything that follows is exploratory.
On 200 episodes four models give the picture of Figure~\ref{fig:weather}.

\explainedfigure{fig_success_by_weather_calibrated_perception_200_episodes}{0.97}{Success rate in four weather conditions under the measured perception degradation}{Four models were each evaluated on the same 200 episodes in clear weather, fog of intensity 1.0, rain of intensity 0.6 and rain of intensity 1.0. In every weather condition the input of the model was degraded as measured: vehicles can disappear from the bird's-eye-view grid, or be placed at a wrong distance. Black lines are 95\% intervals.}{The proposed model stays between 0.89 and 0.91 in all four conditions; the weather-degradation-trained model between 0.88 and 0.90. The reinforcement-learning-fine-tuned model is lower in every condition (0.765 to 0.815) and the baseline trained with traffic is flat at about 0.75.}{Under the perception degradation that was measured and implemented, weather does not lower driving success in sparse traffic. It follows that a robustness or weather-adaptation claim cannot be made from this experiment, and that a weather-adaptive method would have nothing to adapt to here. The next research step is to find the conditions (denser traffic, harsher loss of vehicles) under which weather begins to matter.}{fig:weather}

\begin{table}[H]\centering\small
\caption{Full numbers for Figure~\ref{fig:weather}.}\label{tab:weather}
\resizebox{\textwidth}{!}{\begin{tabular}{llrrrrrr}\toprule
Model & Weather & Success & 95\% interval & Collision & Left road & Timeout & Energy (Wh per km) \\\midrule
Proposed Bird's-Eye-View Driving Model & clear 0.0 & 0.900 & [0.851, 0.934] & 0.075 & 0.020 & 0.005 & 97.9 \\
Proposed Bird's-Eye-View Driving Model & fog 1.0 & 0.890 & [0.839, 0.926] & 0.085 & 0.020 & 0.005 & 96.2 \\
Proposed Bird's-Eye-View Driving Model & rain 0.6 & 0.905 & [0.856, 0.938] & 0.090 & 0.005 & 0.000 & 95.2 \\
Proposed Bird's-Eye-View Driving Model & rain 1.0 & 0.890 & [0.839, 0.926] & 0.105 & 0.005 & 0.000 & 91.9 \\
Weather-Degradation-Trained Driving Model & clear 0.0 & 0.880 & [0.828, 0.918] & 0.110 & 0.000 & 0.010 & 88.8 \\
Weather-Degradation-Trained Driving Model & fog 1.0 & 0.895 & [0.845, 0.930] & 0.105 & 0.000 & 0.000 & 86.8 \\
Weather-Degradation-Trained Driving Model & rain 0.6 & 0.880 & [0.828, 0.918] & 0.120 & 0.000 & 0.000 & 86.9 \\
Weather-Degradation-Trained Driving Model & rain 1.0 & 0.890 & [0.839, 0.926] & 0.110 & 0.000 & 0.000 & 86.2 \\
Reinforcement-Learning-Fine-Tuned Model & clear 0.0 & 0.815 & [0.755, 0.863] & 0.170 & 0.015 & 0.000 & 104.0 \\
Reinforcement-Learning-Fine-Tuned Model & fog 1.0 & 0.780 & [0.718, 0.832] & 0.195 & 0.025 & 0.000 & 104.2 \\
Reinforcement-Learning-Fine-Tuned Model & rain 0.6 & 0.765 & [0.702, 0.818] & 0.215 & 0.020 & 0.000 & 104.2 \\
Reinforcement-Learning-Fine-Tuned Model & rain 1.0 & 0.795 & [0.734, 0.845] & 0.200 & 0.005 & 0.000 & 103.4 \\
Baseline Driving Model Trained With Traffic & clear 0.0 & 0.750 & [0.686, 0.805] & 0.120 & 0.130 & 0.000 & 90.8 \\
Baseline Driving Model Trained With Traffic & fog 1.0 & 0.755 & [0.691, 0.809] & 0.120 & 0.125 & 0.000 & 90.5 \\
Baseline Driving Model Trained With Traffic & rain 0.6 & 0.755 & [0.691, 0.809] & 0.120 & 0.125 & 0.000 & 90.5 \\
Baseline Driving Model Trained With Traffic & rain 1.0 & 0.755 & [0.691, 0.809] & 0.120 & 0.125 & 0.000 & 90.5 \\
\bottomrule\end{tabular}
}
\end{table}

\paragraph{Pre-registered hypotheses, final status.}
\textbf{H0} (the perturbation is real): confirmed, for the detector and for stereo-depth validity (Section~\ref{sec:perception}).
\textbf{H1} (it hurts a clean-trained policy): \emph{not} confirmed.
\textbf{H2} (a weather-input policy beats its non-adaptive control): \emph{not testable}; the reinforcement-learning run with the weather input died after step 0, and the non-adaptive control got worse than its start.
\textbf{H3} (training under degradation helps): \emph{not} confirmed (the degraded-perception student 0.888 pooled over the three degraded cells against 0.895 for the proposed model).

\paragraph{Why weather does not matter here (explanations, partly tested).}
(i)~What is lost is almost only the \emph{detection} of vehicles (Section~\ref{sec:perception}); the map, the road and the car's own state are unaffected. (ii)~In sparse traffic few episodes involve a vehicle near the route in the first place; the failures that exist (Section~\ref{sec:failure}) are dominated by sudden crossings and oncoming vehicles that appear late anyway.
(iii)~Weather has no dynamics effect in this setup, so the road remains as grippy as in clear weather. Explanations (ii) and (iii) were not tested by a dedicated experiment.

\subsection{Caution coupled to perception reliability: learned adaptive behaviour (pilot at 40 episodes, confirmed at 200)}
The caution study tested whether a network that receives a caution value derived from the measured reliability of the detector ($c=1-\text{retention}$) drives more slowly and economically when perception is poor.
\begin{table}[H]\centering\small
\caption{Caution study, sparse traffic, seeds 1010--1049 (40 episodes per cell, pilot).}\label{tab:caution}
\resizebox{\textwidth}{!}{\begin{tabular}{@{}llrrrr@{}}\toprule
Model & Weather & Success & Timeout & Speed (m/s) & Energy (Wh per km)\\\midrule
Headline model & clear 0.0 & 0.925 & 0.000 & 7.94 & 93.0\\
Headline model & fog 1.0 & 0.925 & 0.000 & 7.95 & 94.3\\
Headline model & rain 0.6 & 0.900 & 0.000 & 7.95 & 93.6\\
Headline model & rain 1.0 & 0.925 & 0.000 & 7.97 & 89.4\\
Control: caution input fixed to 0 & clear 0.0 & 0.975 & 0.000 & 8.02 & 74.7\\
Control: caution input fixed to 0 & fog 1.0 & 0.975 & 0.000 & 8.02 & 74.6\\
Control: caution input fixed to 0 & rain 0.6 & 0.975 & 0.000 & 8.02 & 74.3\\
Control: caution input fixed to 0 & rain 1.0 & 0.975 & 0.000 & 8.02 & 74.5\\
Caution-conditioned network, caution 0 & clear 0.0 & 0.900 & 0.000 & 8.02 & 84.7\\
Caution-conditioned network, caution 0 & fog 1.0 & 0.900 & 0.000 & 8.02 & 84.5\\
Caution-conditioned network, caution 0 & rain 0.6 & 0.900 & 0.000 & 8.02 & 84.1\\
Caution-conditioned network, caution 0 & rain 1.0 & 0.900 & 0.000 & 8.02 & 84.3\\
Caution-conditioned network, caution from weather & clear 0.0 & 0.900 & 0.000 & 8.02 & 84.7\\
Caution-conditioned network, caution from weather & fog 1.0 & 0.925 & 0.075 & 6.17 & 59.1\\
Caution-conditioned network, caution from weather & rain 0.6 & 0.900 & 0.100 & 5.89 & 57.0\\
Caution-conditioned network, caution from weather & rain 1.0 & 0.850 & 0.150 & 5.58 & 56.1\\
\bottomrule\end{tabular}
}
\end{table}
Findings: the caution mechanism works as designed (mean speed falls from 8.0 to 5.6--6.2\,m/s and the estimated energy from 84.7 to 56--59\,watt-hours per kilometre as soon as weather is present), at the price of more timeouts (7.5--15\%, because the slower car does not always finish within 800 steps) and, at rain 1.0, a lower success (0.85).
Success is \emph{not} higher than without caution (the fixed-zero control reached 0.975 on these 40 episodes but 0.895 on 200). So the evidence shows \textbf{learned adaptive behaviour}: the student, trained by dataset aggregation from a teacher whose speed and gap depend on the caution value, reproduces that dependence on its own (it slows down in weather; the control with the input fixed at zero does not). It also shows a speed--energy trade-off. It does \emph{not} show a benefit to success: weather does not hurt the proposed model in the first place, so there was no loss for the adaptation to recover. Two limits: what caution means (30 down to 20 km/h, gap 10 up to 25 m) is defined by the teacher, and the caution value is supplied from the known weather (one minus measured detector retention), not sensed by the model from its own perception.

\subsection{Input design (pilot, 10 episodes)}
Five input designs were compared early in the project with one training seed each and 10 episodes per cell, hence wide intervals (Table~\ref{tab:ablation}).
\begin{table}[H]\centering\small
\caption{Input-design ablation, seeds 1000--1009 (10 episodes per cell). Intervals are 95\% Wilson intervals and are very wide.}\label{tab:ablation}
\resizebox{\textwidth}{!}{\begin{tabular}{@{}lrrrrrr@{}}\toprule
Input design & Traffic & Success & 95\% interval & Collision & Left road & Energy (Wh per km)\\\midrule
Overlay grid, three vehicle-state inputs (proposed model) & 0.05 & 0.80 & [0.49, 0.94] & 0.20 & 0.00 & 91\\
Overlay grid, three vehicle-state inputs (proposed model) & 0.2 & 0.60 & [0.31, 0.83] & 0.30 & 0.10 & 154\\
Unmodified grid, six vehicle-state inputs & 0.05 & 0.70 & [0.40, 0.89] & 0.10 & 0.20 & 97\\
Unmodified grid, six vehicle-state inputs & 0.2 & 0.40 & [0.17, 0.69] & 0.30 & 0.20 & 172\\
Overlay grid, six vehicle-state inputs & 0.05 & 0.70 & [0.40, 0.89] & 0.20 & 0.10 & 125\\
Overlay grid, six vehicle-state inputs & 0.2 & 0.30 & [0.11, 0.60] & 0.20 & 0.40 & 201\\
Overlay grid, no vehicle-state inputs & 0.05 & 0.30 & [0.11, 0.60] & 0.30 & 0.40 & 249\\
Overlay grid, no vehicle-state inputs & 0.2 & 0.10 & [0.02, 0.40] & 0.70 & 0.00 & 230\\
No grid, six vehicle-state inputs & 0.05 & 0.70 & [0.40, 0.89] & 0.20 & 0.10 & 83\\
No grid, six vehicle-state inputs & 0.2 & 0.10 & [0.02, 0.40] & 0.80 & 0.10 & 115\\
\bottomrule\end{tabular}
}
\end{table}
What can be said: (a)~the overlay grid with all vehicle-state inputs removed is the worst design (0.30 / 0.10), so the grid alone is not enough when the lane geometry is only in the overlay and the car's speed is unknown;
(b)~vehicle-state inputs alone drive sparse traffic acceptably (0.70) but fail dense traffic (0.10), which is what one expects when no vehicle is visible; (c)~the proposed design is nominally the best in both traffic densities but the intervals overlap completely with three others.
The earlier-stage hypothesis that ``some vehicle-state inputs hurt the clone'' (causal confusion; the early scalar ablation found that removing curvature or lateral offset improved a plain multilayer-perceptron clone in zero traffic) motivated the split design and is consistent with these numbers but is \emph{not} established at this sample size.

\subsection{Robustness to an error in the car's pose against the map (pilot, 40 episodes)}
The overlay depends on knowing where the car is relative to the map. A persistent pose error (random sign; the map is drawn from the wrong pose) was applied with five magnitudes; for scalar-only inputs the same error was added to the lateral and heading vehicle-state inputs.
\explainedfigure{fig_success_versus_localisation_error}{0.97}{Success when the estimate of the car's position on the map is wrong by a constant amount}{The proposed model reads the road from a map drawn into its input, so it needs to know where the car is on the map. Here the position estimate is shifted by a constant lateral distance (horizontal axis, metres; a matching small heading error is added). Success over 40 episodes is shown for three models; bands are 95\% intervals.}{All models are close to their no-error success up to 0.5 metres of error. The proposed model falls to 0.70 at 1 metre and 0.275 at 2 metres; the baseline trained with traffic falls from 0.775 to 0.15.}{The map overlay makes the model accurate when the map alignment is exact and fragile when it is not. A real vehicle would need a position estimate better than about half a metre, or a design that does not depend on one source alone (next figure).}{fig:localisation}

\begin{table}[H]\centering\small
\caption{Numbers for Figure~\ref{fig:localisation}.}\label{tab:localisation}
\resizebox{\textwidth}{!}{\begin{tabular}{@{}lrrrrr@{}}\toprule
Model & Lateral error (m) & Heading error (rad) & Success (40 episodes) & 95\% interval & Left road\\\midrule
Headline model (map-overlay view) & 0 & 0 & 0.925 & [0.80, 0.97] & 0.000\\
Headline model (map-overlay view) & 0.25 & 0.02 & 0.900 & [0.77, 0.96] & 0.025\\
Headline model (map-overlay view) & 0.5 & 0.05 & 0.900 & [0.77, 0.96] & 0.025\\
Headline model (map-overlay view) & 1 & 0.1 & 0.700 & [0.55, 0.82] & 0.250\\
Headline model (map-overlay view) & 2 & 0.2 & 0.275 & [0.16, 0.43] & 0.700\\
Dataset-aggregation student, unmodified view & 0 & 0 & 0.875 & [0.74, 0.95] & 0.025\\
Dataset-aggregation student, unmodified view & 0.25 & 0.02 & 0.750 & [0.60, 0.86] & 0.200\\
Dataset-aggregation student, unmodified view & 0.5 & 0.05 & 0.500 & [0.35, 0.65] & 0.450\\
Dataset-aggregation student, unmodified view & 1 & 0.1 & 0.275 & [0.16, 0.43] & 0.625\\
Dataset-aggregation student, unmodified view & 2 & 0.2 & 0.300 & [0.18, 0.45] & 0.625\\
Baseline behaviour cloning, sparse-traffic demonstrations & 0 & 0 & 0.775 & [0.62, 0.88] & 0.150\\
Baseline behaviour cloning, sparse-traffic demonstrations & 0.25 & 0.02 & 0.600 & [0.45, 0.74] & 0.375\\
Baseline behaviour cloning, sparse-traffic demonstrations & 0.5 & 0.05 & 0.550 & [0.40, 0.69] & 0.375\\
Baseline behaviour cloning, sparse-traffic demonstrations & 1 & 0.1 & 0.225 & [0.12, 0.38] & 0.725\\
Baseline behaviour cloning, sparse-traffic demonstrations & 2 & 0.2 & 0.150 & [0.07, 0.29] & 0.800\\
\bottomrule\end{tabular}
}
\end{table}
Then four further models were trained with random pose error (Table~\ref{tab:variants}) and evaluated on the same error draws (Figure~\ref{fig:posetrained}).
\explainedfigure{fig_success_versus_localisation_error_pose_trained_models}{0.97}{Training with position error, and adding a second lane cue}{The proposed model reads the road from a map drawn into its input, so it needs to know where the car is on the map. Here the position estimate is shifted by a constant lateral distance (horizontal axis, metres; a small heading error grows with it). Success over 40 episodes for the proposed model and four models trained with random position errors; bands are 95\% intervals. Dashed lines: the map is the only source of lane position. Solid green and pink lines: a second, independent and noisy lane measurement (lateral offset and heading) is also given to the model. 40 episodes per point.}{Training with position errors alone does not help (dashed lines, 0.25 to 0.28 at 2 metres of error). Adding a second independent lane cue keeps success between 0.875 and 0.95 for all error sizes up to 2 metres (solid lines).}{Robustness to a wrong position estimate comes from \emph{redundancy} (two sources that can disagree), not from noise augmentation. Caution: the second cue is simulated (simulator truth plus noise), so the result shows that the network can exploit redundancy, not that a real lane detector would be accurate enough.}{fig:posetrained}

\textbf{Interpretation.} Redundancy, not augmentation, is what removes the sensitivity: with one cue a shifted map cannot be recovered from; with two cues the network can learn to trust the one that agrees with the car's motion. The second cue is a noisy copy of simulator truth, so the flat curve shows that the \emph{architecture can use redundancy}, not that a real lane detector would deliver that accuracy.
At the highest second-cue noise level (1.0 of maximum) the conditioned model drops to 0.675--0.70 at 1--2\,m of map error while the unconditioned one stays at 0.85--0.875, so being told the error size did not help consistently (single training seed, 40 episodes: differences of $\pm0.1$ are within noise).
